% from ega1-3-fr.tex
\oldpage[I]{109}
nombre fini quelconque $n$ de $S$-préschémas, l'existence de ces produits
résulte de (\hyperref[I.3.2.6-fr]{3.2.6}) par récurrence sur $n$, en remarquant
que
$(X_1\times_S X_2\times_S\cdots\times_S X_{n-1})\times_S X_n$ satisfait à
la définition du produit. L'unicité du produit entraîne, comme dans toute
catégorie, ses propriétés de \emph{commutativité} et d'\emph{associativité}.
Si, par exemple, $p_1$, $p_2$, $p_3$ désignent les projections de
$X_1\times_S X_2\times_S X_3$, et si on identifie ce préschéma à
$(X_1\times_S X_2)\times_S X_3$, la projection dans $X_1\times_S X_2$ est
identifiée à $(p_1,p_2)_S$.
\end{env}

\begin{env}[3.3.6]
\label{I.3.3.6-fr}
Soient $S$, $S'$ deux préschémas, $\varphi:S'\to S$ un morphisme, qui fait de
$S'$ un $S$-préschéma. Pour tout $S$-préschéma $X$, considérons le produit
$X\times_S S'$, et soient $p$ et $\pi'$ ses projections dans $X$ et $S'$,
respectivement. Muni de $\pi'$, ce produit est un $S'$-préschéma~; quand on le
considère comme tel, on le désigne par $X_{(S')}$ ou $X_{(\varphi)}$, et on dit
que c'est le préschéma obtenu par \emph{extension du préschéma de base} de $S$
à $S'$, au moyen du morphisme $\varphi$, ou l'\emph{image réciproque} de $X$
par $\varphi$. On notera que si $\pi$ est le morphisme structural de $X$,
$\theta$ le morphisme structural de $X\times_S S'$, considéré comme
$S$-préschéma, le diagramme
\phantomsection\label{I.3.3.6.base-change-diagram-fr}
\[
  \xymatrix{
    X\ar[d]_{\pi} &
    X_{(S')}\ar[l]_{p}\ar[ld]_{\theta}\ar[d]^{\pi'}\\
    S & S'\ar[l]^{\varphi}
  }
\]
est commutatif.
\end{env}

\begin{env}[3.3.7]
\label{I.3.3.7-fr}
Avec les notations de (\hyperref[I.3.3.6-fr]{3.3.6}), pour tout
$S$-morphisme $f:X\to Y$, on note encore $f_{(S')}$ le $S'$-morphisme
$f\times_S 1:X_{(S')}\to Y_{(S')}$, et on dit que $f_{(S')}$ est
l'\emph{image réciproque} du morphisme $f$ par $\varphi$. $X_{(S')}$ est donc
un \emph{foncteur covariant} en $X$, de la catégorie des $S$-préschémas dans
celle des $S'$-préschémas.
\end{env}

\begin{env}[3.3.8]
\label{I.3.3.8-fr}
Le préschéma $X_{(S')}$ peut encore être considéré comme solution d'un
\emph{problème d'application universelle}~: tout $S'$-préschéma $T$ est aussi
un $S$-préschéma au moyen de $\varphi$~; tout $S$-morphisme $g:T\to X$
s'écrit alors d'une seule manière $g=p\circ f$, où $f$ est un
$S'$-morphisme $T\to X_{(S')}$, comme il résulte de la définition du produit
appliquée aux $S$-morphismes $f$ et $\psi:T\to S'$ (morphisme structural de
$T$).
\end{env}

\begin{proposition}[3.3.9]
\label{I.3.3.9-fr}
\textup{(«~transitivité de l'extension des préschémas de base~»).} Soient
$S''$ un préschéma, $\varphi':S''\to S'$ un morphisme. Pour tout
$S$-préschéma $X$, il existe un isomorphisme canonique fonctoriel du
$S''$-préschéma $(X_{(\varphi)})_{(\varphi')}$ sur le $S''$-préschéma
$X_{(\varphi\circ\varphi')}$.
\end{proposition}

En effet, soient $T$ un $S''$-préschéma, $\psi$ son morphisme structural,
$g$ un $S$-morphisme de $T$ dans $X$ ($T$ étant considéré comme
$S$-préschéma de morphisme structural $\varphi\circ\varphi'\circ\psi$).
Comme $T$ est aussi un $S'$-préschéma de morphisme structural
$\varphi'\circ\psi$, on peut écrire $g=p\circ g'$, où $g'$ est un
$S'$-morphisme $T\to X_{(\varphi)}$, puis $g'=p'\circ g''$, où $g''$ est un
$S''$-morphisme $T\to(X_{(\varphi)})_{(\varphi')}$~:
\phantomsection\label{I.3.3.9.transitivity-diagram-fr}
\[
  \xymatrix{
    X\ar[d]_{\pi} &
    X_{(\varphi)}\ar[l]_{p}\ar[d]_{\pi'} &
    (X_{(\varphi)})_{(\varphi')}\ar[l]_{p'}\ar[d]_{\pi''}\\
    S & S'\ar[l]^{\varphi} & S''\ar[l]^{\varphi'}
  }
\]
